% Independent analytical audit.  This is an input fragment, not a main file.
\section{Exact null-halfspace gauge counterexamples}
\label{sec:gauge-counterexamples}

These calculations concern the conformal-harmonic lapse candidate, including
its constant Bona--Masso variant.  They do not establish an instability of the
original reference lapse, whose pole coefficients are different.  They are
exact for the null-halfspace Minkowski model specified below.  Transferring a
boundary-model obstruction to the original CMC ball requires the localization
argument described at the end of this section.

\subsection{Exact flat background and a constraint-free scalar sector}

Let $x>0$ be the inward normal coordinate and set
\[
 \bar g=dx^2+2\,dt\,dx+dy^2+dz^2,\qquad \Omega=x,
 \qquad \alpha_*=1,\quad \beta_*^x=1.
\]
The inverse conformal metric has $\bar g^{tt}=-1$,
$\bar g^{tx}=1$, and $\bar g^{xx}=0$.  The conformal gradient of $\Omega$
is null and its Hessian vanishes.  The conformal transformation formula
therefore shows that $g_{\rm phys}=x^{-2}\bar g$ is an exact flat vacuum metric.
Its physical mean curvature is $K_*=-3$, whereas $\bar K_*=0$.

Consider the physical infinitesimal diffeomorphism
\[
 \xi^t=-x\psi,\qquad \xi^x=x\psi,\qquad \xi^A=0.
\]
Keeping the coordinate conformal factor $\Omega=x$ fixed gives
\[
 \delta\bar g=-\mathcal L_\xi\bar g+2\psi\bar g.
\]
Write $D=\partial_t-\partial_x$, $G=dx$, and use the flat spatial metric
for all spatial derivatives and trace-free projections.  Direct ADM
decomposition gives
\begin{align}
 h_{ij}&=0,& c&=-2\psi,& u&=xD\psi,&
 v_i&=-\partial_i(x\psi),\label{eq:halfspace-gauge-fields}\\
 a_{ij}&=-\big[x\partial_i\partial_j\psi+
 G_i\partial_j\psi+G_j\partial_i\psi\big]^{\rm TF},
 & k&=-x^2\Delta\psi+x\psi_x+3\psi.
 \label{eq:halfspace-gauge-curvature}
\end{align}
Here $h$ is the unit-determinant spatial metric perturbation,
$c=\delta\chi$, $u=\delta\bar\alpha$, $v=\delta\beta$,
$a=\delta\bar A$, and $k=\delta K_{\rm phys}$.  Set
$\Theta=Z_i=0$, so the physical $\hat K$ perturbation $q$ equals $k$.

These fields solve the linear vacuum Einstein equations because they are a
physical diffeomorphism of an exact flat metric.  The Hamiltonian and momentum
constraints and the spacetime Weyl perturbation vanish identically.  They also
satisfy the metric-derived Dirac freezing condition: with
\[
 (L_{\rm conf}v)_{ij}=\partial_i v_j+\partial_j v_i
 -\frac23\delta_{ij}\partial_kv^k,
\]
one has $L_{\rm conf}v=2a$, hence $\partial_t h=0$ and
$V_i=\partial_jh_{ij}=0$.  The corresponding elliptic shift equation is
$(\Delta v_i+\partial_i\operatorname{div}v/3)=2\partial^ja_{ij}$.

For a regular $\psi$, the boundary values are
\[
 u|_{x=0}=0,\qquad v_A|_{x=0}=0,\qquad v_x|_{x=0}=-\psi.
\]
The inward-normal linear nullness relation is
$v_x=u-(h_{xx}-c)/2$ and is satisfied.  Thus these are genuine coupled
metric/lapse/elliptic-shift perturbations, not a prescribed shift trace treated
in isolation.

The normal conformal-factor perturbation and conformal trace are
\[
 \delta w=u-v_x=x\psi_t+\psi,\qquad
 \delta\bar K_{\hat{}}=\frac{q-3\delta w}{x}
 =-x\Delta\psi+\psi_x-3\psi_t.
\]
Consequently a lapse with constant conformal Bona--Masso coefficient $f>0$,
\[
 Du=-f\,\delta\bar K_{\hat{}},
\]
is equivalent in this sector to
\begin{equation}
 x\big(D^2\psi-f\Delta\psi\big)
 -(1+3f)\psi_t+(1+f)\psi_x=0.
 \label{eq:halfspace-general-f}
\end{equation}
Equivalently,
$D^2\psi-f\Delta\psi=[(1+3f)D\psi+2f\psi_x]/x$.

\subsection{Conformal-harmonic $f=1$: unbounded tangential-frequency growth}

For $f=1$, take a tangential Fourier mode
$\psi=e^{st+iky}F(kx)$, with $k>0$, $s=\sigma k$, and $X=kx$.
Equation~\eqref{eq:halfspace-general-f} becomes
\begin{equation}
 (2\sigma X-2)F'(X)
 +\big[4\sigma-(\sigma^2+1)X\big]F(X)=0.
 \label{eq:harmonic-turning-ode}
\end{equation}
For any positive integer $m$, choose
\[
 \sigma=(2m+3)^{-1/2},\qquad
 F(X)=\exp\!\left(\frac{\sigma^2+1}{2\sigma}X\right)
 (1-\sigma X)_+^m.
\]
For $0<X<1/\sigma$ its logarithmic derivative is
\[
 \frac{F'}F=\frac{\sigma^2+1}{2\sigma}
 -\frac{m\sigma}{1-\sigma X}.
\]
Substitution in~\eqref{eq:harmonic-turning-ode} leaves exactly
$\sigma(2m+3)-1/\sigma=0$.  The function is zero for
$X\geq1/\sigma$ and is $C^{m-1}$ across the turning point.  Its cutoff
therefore introduces no distributional error in the first-order equation;
choosing $m$ sufficiently large also makes all metric and curvature equations
classical to any prescribed finite order.  In particular $F(0)=1$ and
$F'(0)=2\sigma$.  Every boundary condition stated above is satisfied for all
times.

On a tangential torus, normalize the initial amplitude in any fixed finite
Sobolev norm containing the required metric, lapse, and curvature fields.
Choose $m$ once, larger than the derivative count required by that norm.
The same state at time $t$ is multiplied by $e^{\sigma kt}$.  Thus no
frequency-uniform estimate with a growth constant independent of $k$ can hold
for this exact halfspace model in that finite Sobolev class.  Constraint or
Weyl estimates cannot remove the example: all those fields are identically
zero.  Requiring analyticity or imposing an additional gauge domain would
change the problem; neither is part of the proposed free hybrid.

\subsection{Strictly subluminal $f=1/3$: compatible data with singular outflow}

Absence of a decaying growing normal mode for $0<f<1$ is not sufficient to
prove preservation of regular scri.  At the particular subluminal value
$f=1/3$, equation~\eqref{eq:halfspace-general-f} has the exact solutions
\begin{equation}
 \psi(t,x)=\frac{F(t+x/v)}{x},\qquad
 v=1\pm\frac1{\sqrt3},
 \label{eq:subluminal-singular-outflow}
\end{equation}
for arbitrary smooth $F$.  To check this, set
$\lambda=1/v$ and substitute $x^{-1}F(t+\lambda x)$.
The coefficients of $F/x^2$ and $F'/x$ cancel identically.  The coefficient
of $F''$ is $1-2\lambda+(2/3)\lambda^2=0$.

Choose $F\in C^\infty_c((T_0,T_1))$ with $T_0>0$.
Both $\psi(0,x)$ and $\psi_t(0,x)$ are smooth and compactly supported in
$[vT_0,vT_1]$, away from scri.  Every boundary compatibility jet initially
vanishes.  The solution remains zero near scri until the outgoing front
arrives.  At a later time for which $F(t)\ne0$,
$\psi\sim F(t)/x$, so $c=-2\psi$ and the conformal lapse perturbation
are unbounded.  The constraints and Weyl perturbation remain zero in the
interior.

This example is a regularity obstruction, not an all-time regular solution of
the mixed boundary problem: the singular continuation has
$\xi^x=x\psi\to F(t)$ and therefore ceases to be scri-fixing.
In general nullness also fails once the front reaches the boundary.
A proof of nonexistence for the complete elliptic boundary-coupled problem
must address the possible $S\ne0$ spatial York response.  A bounded smooth
York potential cannot cancel the $1/x$ metric trace; in the tangentially
constant halfspace sector its normal component satisfies $S_x''=0$.
The counterexample nonetheless rules out recommending subluminal slicing
merely on a characteristic-count or decaying-normal-mode test.

\subsection{Relation to the original CMC ball}

The halfspace background is an exact flat physical geometry and is the leading
normal model obtained by magnifying the compactified CMC ball near scri.
It is not the exact global CMC reference, whose conformal lapse, shift, and
factor vary with radius.  A full transfer requires tangential localization,
rescaling $x=X/k$, preservation of the elliptic boundary conditions, and a
uniform estimate on the coefficient and localization remainders.  One useful
route is to start with physical diffeomorphisms of the actual Minkowski
reference, preserving zero constraints and zero Weyl exactly, and correct the
metric-derived connection by the elliptic York solve.  The correction must be
estimated rather than assumed negligible.  Bounded reference-source terms
such as those proportional to the reference conformal mean curvature are
lower order in the magnified problem, but this observation alone is not the
completed quasimode argument.

The conclusions justified here are therefore: an exact frequency-uniform
estimate fails for the conformal-harmonic null-halfspace model; the $f=1/3$
model has exact singular outflow from compatible compact initial data; and
neither counterexample is cured by subsidiary constraints or physical Weyl
control.  An exact global CMC-ball instability theorem is not claimed.
